\section{The temporal equilibrium density}
\label{S:temporal-equilibrium}

The density in Lemma~\ref{L:temporal-equilibrium} is the classical interval
Riesz equilibrium density. Its probability-normalized form on $[-1,1]$ is
recalled by Calef and Hardin~\cite{calef.hardin:source2008:riesz}*{p.~1}.
We prove the potential identity with the normalization used here.

We use Euler's function $\Gamma_E(x) = \int_0^\infty u^{x-1}e^{-u}\ud u$, $x >
0$, and $\operatorname{Beta}(x,y) = \int_0^1t^{x-1}(1-t)^{y-1}\ud t$, $x,y > 0$.
Both integrals are positive and finite.

\LeanMarker{L:equilibrium-gamma-beta}
\begin{lemma}[Gamma and Beta identities]\label{L:equilibrium-gamma-beta}
  For positive arguments,
  \begin{align*}
    \Gamma_E(x+1)&=x\Gamma_E(x),\\
    \operatorname{Beta}(x,y)
    &=\frac{\Gamma_E(x)\Gamma_E(y)}{\Gamma_E(x+y)},\\
    \Gamma_E(2x)
    &=\frac{2^{2x-1}\Gamma_E(x)\Gamma_E(x+1/2)}{\sqrt\pi}.
  \end{align*}
\end{lemma}

\begin{proof}
  Integration by parts gives the recurrence, with vanishing boundary terms. In
  the positive product Gamma integral, use $(u,v) = (Rt,R(1-t))$, with Jacobian
  $R$; Tonelli gives the Beta--Gamma identity. For the last identity, symmetry,
  then $t = (1+v)/2$ and $v^2 = w$, give $\operatorname{Beta}(x,x) =
  2^{1-2x}\operatorname{Beta}(1/2,x)$. Also $\Gamma_E(1/2) = \int_\R e^{-z^2}\ud
  z = \sqrt\pi$: the substitution $u = z^2$ gives the first equality, and
  squaring the positive Gaussian integral and using polar coordinates gives
  $2\pi\int_0^\infty re^{-r^2}\ud r = \pi$. Substituting these into the
  Beta--Gamma identity proves duplication. This proves
  Lemma~\ref{L:equilibrium-gamma-beta}.
\end{proof}

\LeanMarker{L:temporal-equilibrium}
\begin{lemma}[Exact potential and mass]\label{L:temporal-equilibrium}
  For $0 < a < 1$ put $r = (1-a)/2$ and, for $0 < t < 1$,
  \begin{equation}\label{E:temporal-equilibrium-density}
    f_a(t)=\frac{[t(1-t)]^{-r}}{\Gamma_E(r)\Gamma_E(1-r)}.
  \end{equation}
  This positive function belongs to $L^1(0,1)\cap L^2(0,1)$, and
  \begin{align}
    \int_0^1|s-t|^{-a}f_a(t)\ud t&=1
    &&\text{for every }s\in[0,1],\label{E:temporal-equilibrium-potential}\\
    M_a:=\int_0^1f_a(t)\ud t
    &=\frac{\Gamma_E((1+a)/2)}
    {\Gamma_E(1+a)\Gamma_E((1-a)/2)}>0.
    \label{E:temporal-equilibrium-mass}
  \end{align}
\end{lemma}

\begin{proof}
  We first justify the positive series used in the integral calculation. We
  compute the two halves of the potential through an incomplete Beta integral,
  treat the endpoints directly, and then compute the mass and integrability.


  \medskip
  \noindent\textit{Step 1.~}We justify the positive binomial expansion.
  The positive binomial series $(1-z)^{-p} = \sum_{n\ge0}(p)_nz^n/n!$, $p > 0$,
  $0 \le z < 1$, follows directly as follows. Its coefficient ratio gives
  uniform convergence on compact subintervals, including for its derivative.
  Termwise differentiation gives $(1-z)y' = py$, $y(0) = 1$; hence $(1-z)^py$ is
  constant one. Here $(p)_0 = 1$, $(p)_n = p(p+1)\cdots(p+n-1)$. All terms used
  below are nonnegative, so Tonelli applies.


  \medskip
  \noindent\textit{Step 2.~}We compute the potential at interior points and
  endpoints.
  For $0 < s < 1$ split the integral of $|s-t|^{2r-1}[t(1-t)]^{-r}$ into left
  and right parts. In the left part, $t = su$ gives
  \begin{align*}
    L(s)
    &=s^r\int_0^1(1-u)^{2r-1}u^{-r}(1-su)^{-r}\ud u\\
    &=\operatorname{Beta}(1-r,2r)s^r
    \sum_{n\ge0}
    \frac{(r)_n(1-r)_ns^n}{(1+r)_n n!}.
  \end{align*}
  The term of order $n$ integrates $u^{n-r}(1-u)^{2r-1}$, giving
  $\operatorname{Beta}(n+1-r,2r)$. The Gamma recurrence divides this by
  $\operatorname{Beta}(1-r,2r)$ to give $(1-r)_n/(1+r)_n$. Thus, the displayed
  series follows term by term from the positive binomial expansion.
  The Beta ratio follows from the Gamma recurrence. Since $(r)_n/(1+r)_n =
  r/(r+n)$, expanding the incomplete Beta integral gives
  \[
    L(s)=r\operatorname{Beta}(1-r,2r)B_s(r,r),\qquad
    B_s(r,r)=\int_0^s u^{r-1}(1-u)^{r-1}\ud u.
  \]
  Expanding $(1-u)^{r-1}$ in $B_s(r,r)$ gives coefficient $(1-r)_n/n!$ and
  integral $s^{r+n}/(r+n)$. Multiplication by $r\operatorname{Beta}(1-r,2r)$
  therefore gives exactly the coefficients of $L(s)$. Reflection in the time
  interval then supplies the other half of the potential.
  Reflection makes the right part $L(1-s)$, while symmetry gives
  $B_s(r,r)+B_{1-s}(r,r) = \operatorname{Beta}(r,r)$. The full integral is
  therefore $r\operatorname{Beta}(1-r,2r)\operatorname{Beta}(r,r) =
  \Gamma_E(r)\Gamma_E(1-r)$. All integrals are finite because $-r > -1$ and
  $2r-1 > -1$. At $s = 0$ the integral is directly $\operatorname{Beta}(r,1-r) =
  \Gamma_E(r)\Gamma_E(1-r)$; $s = 1$ follows by reflection. This proves
  \eqref{E:temporal-equilibrium-potential}, including the endpoints without a
  limiting argument there.


  \medskip
  \noindent\textit{Step 3.~}We verify both integrability classes and the
  stated mass.
  Since $0<a<1$, one has $0<r<1/2$. Near either endpoint the density behaves
  as a constant times the distance to that endpoint to the power $-r$; its
  square has power $-2r$. Both powers exceed $-1$, so both endpoint integrals
  are finite.
  Integrability and square integrability follow from $-r > -1$ and $-2r = a-1 >
  -1$. Finally,
  \[
    M_a=
    \frac{\operatorname{Beta}(1-r,1-r)}
    {\Gamma_E(r)\Gamma_E(1-r)}
    =\frac{\Gamma_E(1-r)}{\Gamma_E(2-2r)\Gamma_E(r)},
  \]
  which is~\eqref{E:temporal-equilibrium-mass}. Its time mass has not been
  renormalized. This proves Lemma~\ref{L:temporal-equilibrium}.
\end{proof}
